% Generated draft: requires mathematical and source review.
\documentclass[11pt]{article}
\usepackage[margin=25mm]{geometry}
\usepackage{fontspec}
\setmainfont{Latin Modern Roman}
\usepackage{amsmath,amssymb,mathtools,unicode-math}
\setmathfont{Latin Modern Math}
\usepackage{enumitem,needspace,xurl,hyperref}
\setlength{\parindent}{0pt}
\setlength{\parskip}{5pt}
\setlength{\emergencystretch}{3em}
\newcommand{\N}{\mathbb{N}}
\newcommand{\Z}{\mathbb{Z}}
\newcommand{\Q}{\mathbb{Q}}
\newcommand{\R}{\mathbb{R}}
\newcommand{\C}{\mathbb{C}}
\newcommand{\F}{\mathbb{F}}
\newcommand{\A}{\mathbb{A}}
\newcommand{\PP}{\mathbb{P}}
\newcommand{\E}{\mathbb{E}}
\newcommand{\eps}{\varepsilon}
\newcommand{\dd}{\,\mathrm{d}}
\begin{document}
\section*{UnsolvedMath ID: 20001101}
\textbf{Source problem code:} catalogue entry
\subsection{Definitions and mathematical statement}
Let \(G\) be any abelian group, written additively, and let \(A,B\subset G\) be finite with
\[
 A\cap(-B)=\{0\}.
\]
Thus both contain zero, and \(a+b=0\), \(a\in A,b\in B\), forces \(a=b=0\). Define
\[
 A+B=\{a+b:a\in A,b\in B\},\qquad
 A\mathbin{\dot+}B=\{a+b:a\in A,b\in B,a\ne b\}.
\]
V. Lev's source conjecture is
\[
 |A\mathbin{\dot+}B|\ge|A|+|B|-3.
\]
It is motivated by Peter Scherk's 1955 theorem, answering Leo Moser's problem, that the same unique-zero hypothesis gives \(|A+B|\ge|A|+|B|-1\). The source's sharpness examples are \(A=\{0,d,\ldots,(m-1)d\}\), \(B=\{0,d,\ldots,(n-1)d\}\), for positive \(m,n\) and \(d\) of order at least \(m+n-1\), including infinite order.

The original remark reduces the conjecture to nested sets using \(A^*=A\cup B\), \(B^*=A\cap B\). Their total cardinality is unchanged, they retain the unique-zero hypothesis, and their restricted sumset is contained in \(A\mathbin{\dot+}B\). It reports computational verification for all cyclic groups of order at most \(25\), and for \(A=B\) through cyclic order \(36\), plus proofs for torsion-free abelian groups, cyclic prime-order groups, and elementary abelian \(2\)-groups. These are supplied historical checks and special-case claims, not newly verified exhaustive computations or a current literature assessment.

The aid's exact deficit identity holds for arbitrary finite \(A,B\), even without the unique-zero condition. Put \(I=A\cap B\), define the doubling image
\[
 2\cdot I=\{2x:x\in I\},
\]
which is not the sumset \(I+I\), and let
\[
 D=(A+B)\setminus(A\mathbin{\dot+}B)
 =\{s\in A+B:\text{every representation }s=a+b\text{ has }a=b\}.
\]
The reported statements are
\[
 D\subseteq2\cdot I,\qquad
 |A\mathbin{\dot+}B|=|A+B|-|D|.
\]
Consequently, under the source's unique-zero hypothesis,
\[
 |A\mathbin{\dot+}B|\ge|A|+|B|-1-|2\cdot I|.
\]
This proves the requested bound whenever \(|2\cdot I|\le2\), equivalently the overlap meets at most two fibers of the doubling homomorphism, or at most two cosets of \(G[2]=\{g:2g=0\}\). The aid reports sharpness of this special case; the source progression family with \(A=B=\{0,d\}\) and \(d\) of order at least three has restricted sumset \(\{d\}\), attaining \(|A|+|B|-3=1\).

For each \(s\in D\), choose one \(x_s\in I\) with \(2x_s=s\). The aid reports that this selector has no equation \(x+z=2y\) with distinct endpoints \(x,z\); injectivity of doubling on the selector also rules out \(x=z\ne y\). Thus no nonconstant three-term progression is contained in the selector. This is a necessary structural obstruction to a larger pure-diagonal deficit, not a proof of the full conjecture when the overlap has more than two doubling values.
\subsection{Short English statement}
Lev's conjecture asks for a restricted-sumset analogue of Scherk's unique-zero theorem. The supplied aid identifies exactly which ordinary sums are lost by excluding equal summands and proves the desired bound when the overlap has at most two distinct doubles, while retaining the original special cases and computational qualifications.
\subsection{Sources}
\begin{itemize}
\item Local UnsolvedMath registry, permanent ID \(20001101\), source code \texttt{catalogue entry}; \texttt{aim-combinatorics-notes.json}, \texttt{notes[225]}, tag \texttt{problem}.
\item AIM workshop ``Recent trends in additive combinatorics,'' Problem 2.1, V. Lev: \url{https://aimath.org/WWN/additivecomb/additivecomb.pdf}. Statement status \texttt{reconstructed\_unverified}, with extraction artifacts including the inequality \(a\ne b\). No PDF was accessed for this draft.
\item Peter Scherk, 1955 theorem solving a problem by Leo Moser, and the original remark's computational and group-class special cases. No full bibliographic details or computation artifacts are supplied.
\item Research attempt 1: \texttt{partial\_result}, problem status at run \texttt{open}, rigor label \texttt{proved}, novelty confidence low; registry status \texttt{partially\_solved}, classification \texttt{PARTIAL-PROGRESS}. Created and updated 14 August 2026. No original literature field or authoritative external citation for the deficit lemma is supplied. This machine-generated aid is not peer reviewed; independently verify before citation.
\item Ulam AI / UnsolvedMath Contributors; CC BY 4.0.
\end{itemize}
\end{document}
